\documentclass[12pt,a4paper]{article}
% Preambula terpadu untuk tujuh akar mandiri catatan teori representasi Duncan.
% Berkas akar terjemahan tidak diubah; docmute mengabaikan preambula lokalnya.
\usepackage[indonesian]{babel}
\usepackage{fontspec}
\defaultfontfeatures{Ligatures=TeX,Scale=MatchLowercase}
\setmainfont{TeX Gyre Termes}
\setsansfont{TeX Gyre Heros}
\setmonofont{Latin Modern Mono}

\usepackage{amssymb,amsmath,amsthm}
\usepackage{ytableau}
\usepackage{geometry}
\geometry{a4paper,margin=28mm,headheight=16pt,footskip=30pt}
\usepackage{microtype}
\usepackage{csquotes}
\DeclareQuoteAlias{english}{indonesian}
\usepackage{docmute}
\usepackage{fancyhdr}
\usepackage{lastpage}
\usepackage[
  backend=biber,
  style=alphabetic,
  sorting=nyt,
  giveninits=true,
  maxbibnames=99,
  maxcitenames=3
]{biblatex}
\DeclareLanguageMapping{indonesian}{english}
\addbibresource{../source/rep_theory.bib}
\usepackage{xurl}
\usepackage[
  unicode,
  colorlinks=false,
  pdfborder={0 0 0},
  pdftitle={Catatan Teori Representasi - adaptasi bahasa Indonesia},
  pdfauthor={Alexander Duncan},
  pdfsubject={Adaptasi bahasa Indonesia dari catatan MATH 742},
  pdfkeywords={teori representasi, aljabar, grup hingga, CC BY 4.0},
  pdfcreator={OpenAI Codex gpt-5.6-sol, Ultra}
]{hyperref}

% Gaya teorema terpadu. Semua lingkungan bernomor memakai pencacah yang sama.
\theoremstyle{plain}
\newtheorem{theorem}{Teorema}[section]
\newtheorem{lemma}[theorem]{Lema}
\newtheorem{proposition}[theorem]{Proposisi}
\newtheorem{corollary}[theorem]{Korolari}
\newtheorem{conjecture}[theorem]{Konjektur}
\newtheorem{exercise}[theorem]{Latihan}
\theoremstyle{definition}
\newtheorem{definition}[theorem]{Definisi}
\newtheorem{example}[theorem]{Contoh}
\theoremstyle{remark}
\newtheorem{remark}[theorem]{Catatan}
\numberwithin{equation}{section}

% Identitas edisi dan provenance yang tidak berubah antarbangunan.
\newcommand{\DuncanSourceAuthor}{Alexander Duncan}
\newcommand{\DuncanSourceCourse}{MATH 742, Musim Semi 2023}
\newcommand{\DuncanSourceCommit}{c62d36f41189da4bd3da4671668f68720df54ff7}
\newcommand{\DuncanSourceTree}{e83ee440666133b14dec440158a108069a13e9e4}
\newcommand{\DuncanSourceRepository}{https://github.com/vtorsor/representation-theory-notes}
\newcommand{\DuncanLicenseName}{Creative Commons Attribution 4.0 International}
\newcommand{\DuncanLicenseSPDX}{CC-BY-4.0}
\newcommand{\DuncanLicenseURL}{https://creativecommons.org/licenses/by/4.0/}
\newcommand{\DuncanModelProvenance}{OpenAI Codex gpt-5.6-sol, Ultra, atas instruksi pengguna}
\newcommand{\DuncanEditionDate}{Edisi sumber: Musim Semi 2023}

\pagestyle{fancy}
\fancyhf{}
\lhead{MATH 742}
\chead{Catatan Teori Representasi}
\rhead{Adaptasi id-ID}
\lfoot{\DuncanEditionDate}
\rfoot{\thepage\ dari \pageref{LastPage}}
\renewcommand{\footrulewidth}{0.4pt}
\setlength{\parindent}{1.5em}
\setlength{\parskip}{0pt}
\setlength{\emergencystretch}{3em}
\date{\DuncanEditionDate}

% \maketitle pada kelas article menonaktifkan dirinya secara global.
% Simpan definisi awal agar setiap akar mandiri tetap memperoleh halaman judul.
\let\DuncanArticleMakeTitle\maketitle


\begin{document}

\begin{titlepage}
\centering
\vspace*{0.12\textheight}
{\Huge\bfseries Catatan Teori Representasi\par}
\vspace{1.5em}
{\LARGE Adaptasi Bahasa Indonesia\par}
\vspace{3em}
{\Large Karya sumber: \DuncanSourceAuthor\par}
\vspace{0.75em}
{\large \DuncanSourceCourse\par}
\vfill
{\large Lisensi materi sumber dan adaptasi: \DuncanLicenseName\
(\DuncanLicenseSPDX)\par}
\vspace{0.75em}
{\normalsize Adaptasi bahasa dan penyuntingan matematis berbantuan
\DuncanModelProvenance.\par}
\vspace*{0.08\textheight}
\end{titlepage}

\section*{Provenance, atribusi, dan batas komponen}
\addcontentsline{toc}{section}{Provenance, atribusi, dan batas komponen}

Karya sumber disusun oleh \DuncanSourceAuthor\ untuk mata kuliah
\DuncanSourceCourse\ di University of South Carolina. Adaptasi ini diturunkan
dari repositori \url{\DuncanSourceRepository}, komit
\nolinkurl{\DuncanSourceCommit}, pohon \nolinkurl{\DuncanSourceTree}.

Materi sumber tersedia menurut \href{\DuncanLicenseURL}{\DuncanLicenseName}
(\DuncanLicenseSPDX). Berkas lisensi lengkap disertakan sebagai
\texttt{DUNCAN-LICENSE-CC-BY-4.0.txt}. Adaptasi ini menandai perubahan bahasa
dan penyuntingan; tidak ada dukungan, persetujuan, atau pengesahan oleh
\DuncanSourceAuthor, University of South Carolina, Creative Commons, atau
penerbit rujukan yang dinyatakan maupun tersirat.

Adaptasi bahasa Indonesia dan penyuntingan matematis disiapkan dengan bantuan
\DuncanModelProvenance. Model tersebut bukan penulis karya sumber dan tidak
menggantikan atribusi kepada \DuncanSourceAuthor.

Cakupan berlisensi yang dipakai tepat terdiri atas tujuh akar \TeX\ mandiri,
bibliografi bersama, README, dan lisensi dalam pohon repositori tersemat.
Enam lembar tugas pada situs penulis (dilaporkan berjumlah 49 soal) serta satu
solusi parsial berada di luar pohon CC BY 4.0 yang disematkan; semuanya
dikecualikan, tidak diunduh, tidak dimasukkan, tidak diadaptasi, dan tidak
dilisensikan ulang dalam penutupan bangunan ini.

\tableofcontents
\clearpage

% Akar-akar lama memakai BibTeX. Driver terpadu memakai satu bibliografi Biber.
\renewcommand{\bibliographystyle}[1]{}
\renewcommand{\bibliography}[1]{}

% Ke-32 label sumber sudah unik di seluruh tujuh akar, sehingga nama label dan
% rujukannya dipertahankan verbatim saat akar dimasukkan.
\newcommand{\DuncanIncludeRoot}[2]{%
  \clearpage
  \phantomsection
  \addcontentsline{toc}{part}{#2}%
  \begingroup
    \setcounter{section}{0}%
    \setcounter{subsection}{0}%
    \setcounter{subsubsection}{0}%
    \setcounter{theorem}{0}%
    \setcounter{equation}{0}%
    \setcounter{figure}{0}%
    \setcounter{table}{0}%
    \setcounter{footnote}{0}%
    \date{\DuncanEditionDate}%
    \global\let\maketitle\DuncanArticleMakeTitle
    \input{#1}%
  \endgroup
}

\DuncanIncludeRoot{../source/rep_intro.tex}{Pengantar Teori Representasi}
\DuncanIncludeRoot{../source/lin_alg.tex}{Aljabar Multilinear}
\DuncanIncludeRoot{../source/modules.tex}{Modul dan Teori Wedderburn}
\DuncanIncludeRoot{../source/characters.tex}{Teori Karakter}
\DuncanIncludeRoot{../source/induction.tex}{Representasi Terinduksi}
\DuncanIncludeRoot{../source/symmetric.tex}{Grup Simetrik dan Grup Linear Umum}
\DuncanIncludeRoot{../source/nonclosed.tex}{Representasi atas Medan yang Tidak Tertutup}

\clearpage
\printbibliography[heading=bibintoc,title={Daftar Pustaka}]

\end{document}
